% AMC12B 2017 Problem 10 - Strategic Analysis
% Using the AMC12/COMC Strategic Problem Analysis Framework
\documentclass[]{article}
\usepackage[utf8]{inputenc}
\usepackage{amsmath, amssymb, amsthm}
\usepackage{geometry}
\usepackage{xcolor}
\usepackage[most]{tcolorbox}
\usepackage{enumitem}
\usepackage{fancyhdr}

% Page setup
\geometry{margin=1in}
\pagestyle{fancy}
\fancyhf{}
\rhead{AMC12B 2017 Problem 10}
\lhead{\today}

% Color definitions
\definecolor{problemconfig}{RGB}{230, 242, 255}
\definecolor{strategic}{RGB}{255, 230, 242}
\definecolor{tactical}{RGB}{242, 255, 230}
\definecolor{techniques}{RGB}{255, 242, 230}
\definecolor{verification}{RGB}{255, 230, 230}
\definecolor{solution}{RGB}{230, 255, 255}
\definecolor{competition}{RGB}{242, 242, 255}
\definecolor{gold}{RGB}{218, 165, 32}

% Box styles
\newtcolorbox{problemconfigbox}{
    colback=problemconfig,
    colframe=blue,
    title=PROBLEM CONFIGURATION ANALYSIS,
    fonttitle=\bfseries,
    arc=3pt,
    boxrule=1pt,
    breakable
}

\newtcolorbox{strategicbox}{
    colback=strategic,
    colframe=purple,
    title=STRATEGIC FRAMEWORK APPLICATION,
    fonttitle=\bfseries,
    arc=3pt,
    boxrule=1pt,
    breakable
}

\newtcolorbox{tacticalbox}{
    colback=tactical,
    colframe=green,
    title=TACTICAL METHODOLOGY SELECTION,
    fonttitle=\bfseries,
    arc=3pt,
    boxrule=1pt,
    breakable
}

\newtcolorbox{techniquesbox}{
    colback=techniques,
    colframe=orange,
    title=CONVERSION TECHNIQUES APPLICATION,
    fonttitle=\bfseries,
    arc=3pt,
    boxrule=1pt,
    breakable
}

\newtcolorbox{verificationbox}{
    colback=verification,
    colframe=red,
    title=VERIFICATION STRATEGY DESIGN,
    fonttitle=\bfseries,
    arc=3pt,
    boxrule=1pt,
    breakable
}

\newtcolorbox{solutionbox}{
    colback=solution,
    colframe=cyan,
    title=SOLUTION ROADMAP,
    fonttitle=\bfseries,
    arc=3pt,
    boxrule=1pt,
    breakable
}

\newtcolorbox{competitionbox}{
    colback=competition,
    colframe=purple,
    title=COMPETITION STRATEGY INTEGRATION,
    fonttitle=\bfseries,
    arc=3pt,
    boxrule=1pt,
    breakable
}

\newtcolorbox{keyinsightbox}{
    colback=yellow!20,
    colframe=gold,
    title=KEY INSIGHT,
    fonttitle=\bfseries,
    arc=3pt,
    boxrule=2pt,
    leftrule=5pt,
    breakable
}

\newtcolorbox{warningbox}{
    colback=red!10,
    colframe=red!50,
    title=WARNING: Common Pitfall,
    fonttitle=\bfseries,
    arc=3pt,
    boxrule=2pt,
    leftrule=5pt,
    breakable
}

\newtcolorbox{tipbox}{
    colback=green!10,
    colframe=green!50,
    title=PRO TIP,
    fonttitle=\bfseries,
    arc=3pt,
    boxrule=2pt,
    leftrule=5pt,
    breakable
}

\begin{document}

\title{\textbf{AMC12B 2017 Problem 10\\Strategic Analysis}}
\author{Competition Math Framework}
\date{\today}
\maketitle

\section*{Problem Statement}

At Typico High School, $60\%$ of the students like dancing, and the rest dislike it. Of those who like dancing, $80\%$ say that they like it, and the rest say that they dislike it. Of those who dislike dancing, $90\%$ say that they dislike it, and the rest say that they like it. What fraction of students who say they dislike dancing actually like it?

\vspace{0.5em}
\textbf{Answer Choices:}\\
$\textbf{(A)}\ 10\%\qquad\textbf{(B)}\ 12\%\qquad\textbf{(C)}\ 20\%\qquad\textbf{(D)}\ 25\%\qquad\textbf{(E)}\ \frac{100}{3}\%$

\begin{problemconfigbox}
\subsection*{A. Problem Classification}
\begin{itemize}[leftmargin=*]
    \item \textbf{Problem Type}: To Find - calculate a conditional probability (fraction)
    \item \textbf{Difficulty Band}: Mid-band (Problem 10 of 25) - requires conditional probability understanding but is accessible with clear reasoning
    \item \textbf{Competition Context}: AMC12 (75 min, 25 problems) - approximately 3 minutes allocated per problem
    \item \textbf{Mathematical Domain}: Probability (conditional probability, Bayes' Theorem)
\end{itemize}

\subsection*{B. Problem Structure Identification}
\begin{itemize}[leftmargin=*]
    \item \textbf{Given Information}:
    \begin{itemize}
        \item $60\%$ of students like dancing (therefore $40\%$ dislike it)
        \item Of those who like dancing: $80\%$ say they like it, $20\%$ say they dislike it
        \item Of those who dislike dancing: $90\%$ say they dislike it, $10\%$ say they like it
    \end{itemize}
    
    \item \textbf{Constraints}:
    \begin{itemize}
        \item Percentages must sum to $100\%$ within each group
        \item Students are partitioned into two groups (like/dislike) with two reporting behaviors (say like/say dislike)
    \end{itemize}
    
    \item \textbf{Objective}: Find $P(\text{Actually Like} \mid \text{Say Dislike})$
    
    \item \textbf{Hidden Assumptions}:
    \begin{itemize}
        \item Students' statements don't always match their true preferences
        \item The problem creates a $2 \times 2$ confusion matrix situation
        \item We need conditional probability, not simple probability
    \end{itemize}
    
    \item \textbf{Answer Choice Leverage}:
    \begin{itemize}
        \item All answers are "nice" percentages or simple fractions
        \item Can eliminate obviously wrong answers through sanity checks
        \item Should be able to compute exact answer without approximation
    \end{itemize}
\end{itemize}

\subsection*{C. Strategic Orientation}
\begin{itemize}[leftmargin=*]
    \item \textbf{Hypothesis}: Given student population with preferences and reporting behaviors
    
    \item \textbf{Conclusion}: Need to find what fraction of those who \textit{say} they dislike actually \textit{like} dancing
    
    \item \textbf{Notation}:
    \begin{itemize}
        \item $L$ = Actually likes dancing ($60\%$ of students)
        \item $D$ = Actually dislikes dancing ($40\%$ of students)
        \item $S_L$ = Says they like dancing
        \item $S_D$ = Says they dislike dancing
    \end{itemize}
    
    \item \textbf{Intuitive Approach}: Use conditional probability formula or construct a $2 \times 2$ table with concrete numbers (assume 100 students for simplicity)
    
    \item \textbf{Cross-Domain Potential}: This is a pure probability problem, but could be visualized with tree diagrams or Venn diagrams
\end{itemize}
\end{problemconfigbox}

\begin{strategicbox}
\subsection*{A. Psychological Strategies}
\begin{itemize}[leftmargin=*]
    \item \textbf{Mental Toughness}: Problem requires careful tracking of percentages - resist the urge to rush through calculations
    
    \item \textbf{Creativity Cultivation}: Consider multiple solution paths (Bayes' Theorem vs. concrete numbers approach)
    
    \item \textbf{Iterative Refinement}: If first approach seems messy, pivot to the concrete numbers method which is often clearer
\end{itemize}

\subsection*{B. Investigation Strategies}
\begin{itemize}[leftmargin=*]
    \item \textbf{Startup Orientation}: Immediately recognize this as a conditional probability problem - the key phrase is "of students who say they dislike"
    
    \item \textbf{Get Your Hands Dirty}: Use concrete numbers (assume 100 students) to avoid percentage confusion
    
    \item \textbf{Penultimate Step}: Need to find two quantities: (1) number who like but say dislike, and (2) total number who say dislike
    
    \item \textbf{Wishful Thinking}: If we had a table showing all four combinations (like/say like, like/say dislike, dislike/say like, dislike/say dislike), we could directly compute the answer
    
    \item \textbf{Make It Easier}: Use 100 students instead of working with abstract percentages
    
    \item \textbf{Visualization}: Create a $2 \times 2$ table or tree diagram to organize the information
\end{itemize}

\subsection*{C. Late-Band Strategic Playbook}
\begin{itemize}[leftmargin=*]
    \item Not applicable - this is a mid-band problem (P10)
    \item However, answer choices can still be used for sanity checking
\end{itemize}

\subsection*{D. Argument Construction Strategy}
\begin{itemize}[leftmargin=*]
    \item \textbf{Direct Proof}: Apply conditional probability formula directly or use population counting
    
    \item \textbf{Method Selection}: Direct calculation with concrete numbers is clearest for this problem type
\end{itemize}

\subsection*{E. Cross-Domain Translation}
\begin{itemize}[leftmargin=*]
    \item \textbf{Draw a Picture}: Tree diagram with two stages: (1) actual preference, (2) stated preference
    
    \item \textbf{Recast the Problem}: This is equivalent to finding the false negative rate in a diagnostic test scenario
    
    \item \textbf{Change Your Point of View}: Can think of this as signal detection theory or confusion matrix analysis
\end{itemize}
\end{strategicbox}

\begin{tacticalbox}
\subsection*{A. Core Mathematical Tactics}
\begin{itemize}[leftmargin=*]
    \item \textbf{Conditional Probability}: The central tactic - computing $P(A \mid B) = \frac{P(A \cap B)}{P(B)}$
    
    \item \textbf{Law of Total Probability}: The denominator requires summing over all ways to "say dislike"
    
    \item \textbf{Partition Principle}: The student population is partitioned into like/dislike groups
\end{itemize}

\subsection*{B. Crossover Tactics}
\begin{itemize}[leftmargin=*]
    \item \textbf{Table/Matrix Organization}: Using a $2 \times 2$ table to organize the four scenarios
    
    \item \textbf{Concrete Value Method}: Using $N = 100$ students to avoid percentage confusion
\end{itemize}
\end{tacticalbox}

\begin{techniquesbox}
\subsection*{A. Recognition Index - Why This Technique Applies}
\begin{itemize}[leftmargin=*]
    \item \textbf{Signature Pattern}: Phrase "of students who say they dislike" signals conditional probability
    
    \item \textbf{Trigger}: Multiple conditional percentages given (80\%, 90\%) suggest Bayes' Theorem or table method
    
    \item \textbf{Problem Structure}: Two-stage process (actual preference $\rightarrow$ stated preference) indicates need for joint and marginal probabilities
\end{itemize}

\subsection*{B. Applicable High-Probability Techniques}
\begin{itemize}[leftmargin=*]
    \item \textbf{Technique 1: Bayes' Theorem (Primary)}
    \begin{itemize}
        \item Formula: $P(L \mid S_D) = \frac{P(S_D \mid L) \cdot P(L)}{P(S_D)}$
        \item $P(S_D \mid L) = 0.2$ (20\% of likers say dislike)
        \item $P(L) = 0.6$ (60\% actually like)
        \item $P(S_D) = P(S_D \mid L) \cdot P(L) + P(S_D \mid D) \cdot P(D)$
        \item $P(S_D) = 0.2 \times 0.6 + 0.9 \times 0.4 = 0.12 + 0.36 = 0.48$
        \item $P(L \mid S_D) = \frac{0.2 \times 0.6}{0.48} = \frac{0.12}{0.48} = \frac{1}{4} = 25\%$
    \end{itemize}
    
    \item \textbf{Technique 2: Concrete Numbers (Alternative)}
    \begin{itemize}
        \item Assume 100 students total
        \item 60 students like dancing: $60 \times 0.8 = 48$ say like, $60 \times 0.2 = 12$ say dislike
        \item 40 students dislike dancing: $40 \times 0.1 = 4$ say like, $40 \times 0.9 = 36$ say dislike
        \item Total who say dislike: $12 + 36 = 48$ students
        \item Of these, 12 actually like dancing
        \item Fraction: $\frac{12}{48} = \frac{1}{4} = 25\%$
    \end{itemize}
\end{itemize}

\subsection*{C. Technique Integration}
Both methods yield the same answer ($25\%$), providing strong confirmation. The concrete numbers method is often clearer for students less familiar with Bayes' Theorem, while Bayes' Theorem is more elegant and generalizable.

\subsection*{D. Rationale for Selection}
\begin{itemize}[leftmargin=*]
    \item \textbf{Why Concrete Numbers}: Eliminates percentage confusion, makes calculation transparent, suitable for competition time pressure
    
    \item \textbf{Why Bayes' Theorem}: More sophisticated, demonstrates understanding of conditional probability, useful for verification
    
    \item \textbf{Competition Choice}: Use concrete numbers as primary method (faster, less error-prone), verify with Bayes if time permits
\end{itemize}
\end{techniquesbox}

\begin{verificationbox}
\subsection*{A. Verification Level Selection - Decision Tree}

Given:
\begin{itemize}[leftmargin=*]
    \item Problem difficulty: Mid-band (P10)
    \item Solution confidence: High (90\%+) after applying both methods
    \item Time available: Approximately 3 minutes allocated
    \item Impact: Standard AMC12 problem, 6 points
\end{itemize}

\textbf{Selected Level}: Level 4 - Standard Verification (30-60 seconds)

\subsection*{B. Specific Verification Methods}

\textbf{Primary Verification:}
\begin{enumerate}[leftmargin=*]
    \item \textbf{Sanity Check}: Does $25\%$ make sense?
    \begin{itemize}
        \item Yes - it's between $0\%$ and $100\%$
        \item 12 out of 48 is indeed exactly $\frac{1}{4}$
        \item Not too small (rules out 10\%, 12\%) nor too large (rules out 33\%)
    \end{itemize}
    
    \item \textbf{Alternative Method}: Verify with Bayes' Theorem
    \begin{itemize}
        \item Computed: $\frac{0.12}{0.48} = \frac{1}{4}$ $\checkmark$
        \item Matches concrete numbers result $\checkmark$
    \end{itemize}
    
    \item \textbf{Boundary Verification}:
    \begin{itemize}
        \item Total students: $48 + 52 = 100$ $\checkmark$
        \item All groups sum correctly $\checkmark$
    \end{itemize}
    
    \item \textbf{Answer Choice Match}:
    \begin{itemize}
        \item $\frac{1}{4} = 25\%$ corresponds to choice (D) $\checkmark$
    \end{itemize}
\end{enumerate}

\textbf{Specialized Domain Verification (Probability):}
\begin{itemize}[leftmargin=*]
    \item \textbf{Probability Axioms}: All probabilities in [0,1] $\checkmark$
    \item \textbf{Conditional Probability Structure}: Numerator $\leq$ Denominator $\checkmark$ ($12 \leq 48$)
    \item \textbf{Consistency}: Creating the full $2 \times 2$ table shows all cells are non-negative $\checkmark$
\end{itemize}

\subsection*{C. Confidence Calibration}

\textbf{Initial Confidence}: 85\% (after concrete numbers method)

\textbf{Post-Verification Confidence}: 95\%
\begin{itemize}[leftmargin=*]
    \item Two independent methods agree
    \item All sanity checks pass
    \item Answer matches a provided choice
    \item No computational errors detected
\end{itemize}

\textbf{Flag for Review}: No - high confidence, verified answer
\end{verificationbox}

\begin{solutionbox}
\subsection*{Step-by-Step Solution Plan}

\textbf{Phase 1: Problem Understanding (30 seconds)}
\begin{enumerate}[leftmargin=*]
    \item Read problem carefully, identifying it as conditional probability
    \item Note the key phrase: "of students who say they dislike"
    \item Recognize we need $P(\text{Actually Like} \mid \text{Say Dislike})$
    \item Time checkpoint: If not clear by 30 sec, re-read
\end{enumerate}

\textbf{Phase 2: Setup (30 seconds)}
\begin{enumerate}[leftmargin=*]
    \item Decide on approach: Choose concrete numbers (assume 100 students)
    \item Identify the two groups and their sizes:
    \begin{itemize}
        \item 60 students actually like dancing
        \item 40 students actually dislike dancing
    \end{itemize}
    \item Time checkpoint: Setup should be quick and clear
\end{enumerate}

\textbf{Phase 3: Calculation (60-90 seconds)}
\begin{enumerate}[leftmargin=*]
    \item Calculate how many students say they dislike from each group:
    \begin{itemize}
        \item From "like" group: $60 \times 0.2 = 12$ students
        \item From "dislike" group: $40 \times 0.9 = 36$ students
    \end{itemize}
    
    \item Find total who say dislike: $12 + 36 = 48$ students
    
    \item Calculate the fraction: $\frac{12}{48} = \frac{1}{4} = 25\%$
    
    \item Time checkpoint: Should have answer by 2 minutes
\end{enumerate}

\textbf{Phase 4: Verification (30-45 seconds)}
\begin{enumerate}[leftmargin=*]
    \item Quick verification with Bayes' Theorem (if time permits)
    \item Check that answer matches choice (D)
    \item Confirm all numbers make sense
    \item Time checkpoint: Total time $\leq$ 3 minutes
\end{enumerate}

\subsection*{Alternative Approaches}

\textbf{Alternative 1: Direct Bayes' Theorem}
\begin{itemize}[leftmargin=*]
    \item Compute $P(L \mid S_D) = \frac{P(S_D \mid L) \cdot P(L)}{P(S_D)}$
    \item More elegant but requires comfort with formula
    \item Use as verification method if time permits
\end{itemize}

\textbf{Alternative 2: Tree Diagram}
\begin{itemize}[leftmargin=*]
    \item Draw tree with two levels
    \item First level: Actually Like (60\%) vs Dislike (40\%)
    \item Second level: Say Like vs Say Dislike
    \item Trace paths that end in "Say Dislike"
    \item More visual but takes more time
\end{itemize}

\subsection*{Pivot Indicators}
\begin{itemize}[leftmargin=*]
    \item If percentages become confusing $\rightarrow$ Switch to concrete numbers
    \item If calculation seems wrong $\rightarrow$ Verify with alternative method
    \item If taking $>$ 4 minutes $\rightarrow$ Mark answer and move on
\end{itemize}

\subsection*{Time Management}
\begin{itemize}[leftmargin=*]
    \item \textbf{Target time}: 2.5-3 minutes (appropriate for P10)
    \item \textbf{Checkpoint at 2 min}: Should have numerical answer
    \item \textbf{Checkpoint at 3 min}: Should have verified and marked answer
    \item \textbf{Abandon threshold}: If no progress by 4 minutes, flag and move on
\end{itemize}
\end{solutionbox}

\begin{competitionbox}
\subsection*{A. Time Management Strategy}
\begin{itemize}[leftmargin=*]
    \item \textbf{Allocated Time}: 3 minutes (appropriate for mid-band problem)
    \item \textbf{Actual Expected Time}: 2.5-3 minutes with verification
    \item \textbf{Time Distribution}:
    \begin{itemize}
        \item Understanding: 30 seconds
        \item Setup: 30 seconds
        \item Calculation: 90 seconds
        \item Verification: 30 seconds
    \end{itemize}
\end{itemize}

\subsection*{B. Problem Prioritization}
\begin{itemize}[leftmargin=*]
    \item \textbf{Solve Now}: Yes - P10 is accessible and worth full points
    \item \textbf{Difficulty Assessment}: Medium - requires conditional probability understanding but straightforward calculation
    \item \textbf{Domain Comfort}: If comfortable with probability, solve immediately; if not, consider flagging for later
\end{itemize}

\subsection*{C. Risk Management}
\begin{itemize}[leftmargin=*]
    \item \textbf{Confidence Threshold}: 95\% after verification - high confidence
    \item \textbf{Answer Choice Validation}: Answer is (D) $25\%$, a clean fraction
    \item \textbf{Error Recovery}: If initial approach confusing, immediately pivot to concrete numbers method
\end{itemize}

\subsection*{D. Abandonment Criteria}
\begin{itemize}[leftmargin=*]
    \item \textbf{Soft Abandon Trigger}: If no clear approach after 1 minute, consider moving to next problem
    \item \textbf{Hard Abandon Trigger}: If calculation becomes messy or confusing after 4 minutes total
    \item \textbf{Return Strategy}: If abandoned, likely indicates need to review conditional probability concepts
\end{itemize}

\subsection*{E. Answer Format Management}
\begin{itemize}[leftmargin=*]
    \item \textbf{Multiple Choice}: Mark (D) carefully on answer sheet
    \item \textbf{Verification}: $25\% = \frac{1}{4}$ matches choice (D) exactly
    \item \textbf{Common Errors to Avoid}:
    \begin{itemize}
        \item Don't bubble (B) 12\% (that's just the numerator as percentage)
        \item Don't bubble (C) 20\% (that's $P(S_D \mid L)$, not what we want)
    \end{itemize}
\end{itemize}
\end{competitionbox}

\begin{keyinsightbox}
\textbf{The Key Realization:}

This problem asks for $P(\text{Actually Like} \mid \text{Say Dislike})$, which requires finding:
\begin{itemize}
    \item \textbf{Numerator}: Number who actually like BUT say dislike (the "false negatives")
    \item \textbf{Denominator}: Total number who say dislike (from BOTH groups)
\end{itemize}

The critical insight is that students who say they dislike come from \textit{two sources}:
\begin{enumerate}
    \item Those who like but say dislike (20\% of the 60\% who like = 12\%)
    \item Those who dislike and say dislike (90\% of the 40\% who dislike = 36\%)
\end{enumerate}

We want the first group as a fraction of the total: $\frac{12}{12+36} = \frac{12}{48} = \frac{1}{4} = 25\%$
\end{keyinsightbox}

\begin{warningbox}
\textbf{Most Common Error: Confusing Conditional Probabilities}

Students often incorrectly compute $P(S_D \mid L) = 20\%$ and select choice (C), but this is NOT what the problem asks for!

\vspace{0.5em}
\textbf{What the problem asks}: Of students who SAY dislike, what fraction ACTUALLY like?

\textbf{What 20\% represents}: Of students who ACTUALLY like, what fraction SAY dislike?

\vspace{0.5em}
These are \textit{not} the same! The problem requires reversing the conditional relationship, which is why we need Bayes' Theorem or the concrete numbers approach.

\vspace{0.5em}
\textbf{Another Common Error:} Computing $\frac{12}{100} = 12\%$ (choice B) by dividing by the total population instead of by those who say dislike. Always check: what is the reference group for "of students who..."?
\end{warningbox}

\begin{tipbox}
\textbf{Competition Strategy for Conditional Probability Problems:}

\begin{enumerate}
    \item \textbf{Identify the conditioning}: Look for "of students who..." phrases
    \item \textbf{Use concrete numbers}: Assume 100 people to eliminate fraction/percentage confusion
    \item \textbf{Build a table}: Create a $2 \times 2$ table showing all four scenarios
    \item \textbf{Verify with formula}: If time permits, check with Bayes' Theorem
\end{enumerate}

\textbf{Time-Saving Tip}: For AMC12 conditional probability problems, the concrete numbers approach is usually faster and less error-prone than manipulating formulas directly. Save Bayes' Theorem for verification!
\end{tipbox}

\section*{Final Answer}
\boxed{\textbf{(D) } 25\%}

\section*{Confidence Assessment}
\begin{itemize}[leftmargin=*]
    \item \textbf{Confidence Level}: 95\% - Two independent methods yield identical answer, all verification checks pass
    
    \item \textbf{Verification Applied}: Level 4 Standard Verification - sanity checks, alternative method (Bayes' Theorem), boundary verification, answer choice match
    
    \item \textbf{Flag for Review}: No - high confidence in both approach and calculation
    
    \item \textbf{Time Spent}: Approximately 2.5-3 minutes (optimal for P10)
\end{itemize}

\section*{Learning Points}

\textbf{Key Takeaways from This Problem:}
\begin{enumerate}
    \item Conditional probability problems require careful attention to which group is the reference (denominator)
    \item Concrete numbers (assuming 100 students) often clarifies percentage problems
    \item Bayes' Theorem provides elegant verification for conditional probability calculations
    \item Creating a $2 \times 2$ table helps organize information and prevents errors
    \item Always distinguish $P(A \mid B)$ from $P(B \mid A)$ - they are generally not equal!
\end{enumerate}

\textbf{Transferable Skills:}
\begin{itemize}
    \item Recognition of conditional probability patterns
    \item Ability to switch between abstract formulas and concrete calculations
    \item Systematic verification using multiple methods
    \item Time management for mid-difficulty problems
\end{itemize}

\end{document}